%% Speaker Notes for LA Kelp Detection Prototype Progress
%% Compile with: pdflatex speaker_notes.tex (or latexmk)

\documentclass[11pt]{article}
\usepackage[utf8]{inputenc}
\usepackage[T1]{fontenc}
\usepackage[margin=1in]{geometry}
\usepackage{amsmath}
\usepackage{booktabs}
\usepackage{xcolor}
\usepackage{enumitem}
\usepackage{hyperref}
\usepackage{titlesec}
\usepackage{listings}

% Code formatting
\lstset{
  basicstyle=\ttfamily\small,
  backgroundcolor=\color{gray!10},
  frame=single,
  framerule=0pt,
  breaklines=true,
  columns=fullflexible
}

% Section formatting
\titleformat{\section}{\large\bfseries\color{blue!70!black}}{}{0em}{}[\titlerule]
\titleformat{\subsection}{\normalsize\bfseries}{}{0em}{}

% Custom environment for talking points
\newenvironment{talkingpoints}{%
  \begin{itemize}[leftmargin=1.5em, itemsep=0.3em]
}{%
  \end{itemize}
}

\title{\textbf{Speaker Notes}\\LA Kelp Detection Prototype\\Weekly Progress Report}
\author{Ekaba Bisong}
\date{\today}

\begin{document}

\maketitle
\tableofcontents
\newpage

%% ============================================================
\section{Slide 1: Title Slide}
%% ============================================================

\begin{talkingpoints}
  \item This presentation covers the progress on building a Learning Automata prototype for kelp detection using Sentinel-2 satellite imagery.
  \item The prototype implements the core ideas from the thesis proposal: using LA to adaptively select detection heuristics.
  \item We'll cover two weeks of work: project setup and initial experiments.
\end{talkingpoints}

%% ============================================================
\section{Slide 2: Outline}
%% ============================================================

\begin{talkingpoints}
  \item Week 1 focused on setup after the committee meeting.
  \item Week 2 focused on implementation and running first experiments.
  \item The key finding: the prototype works, but hyperparameter tuning matters.
\end{talkingpoints}

%% ============================================================
\section{Slide 3: Week 1 -- Committee Meeting \& Project Setup}
%% ============================================================

\subsection{Key Points to Emphasize}

\begin{talkingpoints}
  \item The committee meeting on February 5th provided critical guidance: \textbf{build a full but simple prototype first}.
  \item ``Full'' means end-to-end: data loading, heuristics, LA, evaluation.
  \item ``Simple'' means spectral heuristics only---no deep learning yet.
  \item Important scope clarification: we're doing kelp \textbf{mapping} (binary segmentation), not full ecological monitoring (species identification, health assessment, etc.).
  \item The Global Policy approach uses a single $L_{R-I}$ automaton that learns one probability distribution over heuristics for the entire dataset.
\end{talkingpoints}

\subsection{What is $L_{R-I}$?}

\begin{talkingpoints}
  \item $L_{R-I}$ stands for ``Linear Reward-Inaction.''
  \item \textbf{Reward}: When the selected heuristic performs well, increase its probability and decrease others.
  \item \textbf{Inaction}: When performance is poor, do nothing (don't update probabilities).
  \item This is a conservative scheme---it only reinforces good outcomes.
\end{talkingpoints}

%% ============================================================
\section{Slide 4: Week 1 -- Prototype Architecture}
%% ============================================================

\subsection{The Four Heuristics}

\begin{talkingpoints}
  \item \textbf{NDVI} (Normalized Difference Vegetation Index): Classic vegetation index using NIR and Red bands. Works well for terrestrial vegetation, less tuned for marine.
  \item \textbf{FAI} (Floating Algae Index): Designed specifically for detecting floating vegetation in water. Uses NIR, Red, and SWIR bands.
  \item \textbf{GNDVI} (Green NDVI): Uses Green band instead of Red. Often better for aquatic vegetation.
  \item \textbf{Ensemble}: Takes majority vote across the other three heuristics.
\end{talkingpoints}

\subsection{Infrastructure Decisions}

\begin{talkingpoints}
  \item \texttt{config.yaml}: Single source of truth for all experiment parameters (learning rate, reward type, number of epochs, etc.).
  \item \texttt{Makefile}: Reproducible experiment execution with \texttt{make run} and \texttt{make baselines}.
  \item Experiment tracker auto-logs results, plots, and configuration for each run.
\end{talkingpoints}

%% ============================================================
\section{Slide 5: Week 1 -- Next Steps}
%% ============================================================

\begin{talkingpoints}
  \item The immediate tasks were to build the infrastructure and implement all components.
  \item Success criteria are concrete and measurable:
  \begin{enumerate}
    \item The LA must update probabilities correctly (mechanical correctness).
    \item Probabilities should converge toward better heuristics (learning signal).
    \item LA should beat random selection (practical value).
  \end{enumerate}
  \item If we can't beat random selection, the LA isn't providing value.
\end{talkingpoints}

%% ============================================================
\section{Slide 6: Week 2 -- Prototype Implementation Complete}
%% ============================================================

\begin{talkingpoints}
  \item All core components are now implemented and working.
  \item The data loader handles 10-band Sentinel-2 tiles (coastal aerosol, blue, green, red, vegetation red edge bands, NIR, SWIR).
  \item The reward function supports both binary (threshold-based) and continuous (raw IoU) modes.
  \item We have three baselines for comparison: Random, Fixed, and Oracle.
\end{talkingpoints}

\subsection{What Are the Baselines?}

\begin{talkingpoints}
  \item \textbf{Random}: Select a heuristic uniformly at random for each tile. This is the \textbf{lower bound}---if LA can't beat this, it's useless.
  \item \textbf{Fixed}: Always use the same heuristic for every tile. One fixed baseline per heuristic.
  \item \textbf{Oracle}: Try all heuristics, pick the best one per tile (using ground truth). This is the \textbf{upper bound}---perfect hindsight selection.
\end{talkingpoints}

%% ============================================================
\section{Slide 7: Week 2 -- Baseline Results}
%% ============================================================

\subsection{Understanding the Numbers}

\begin{talkingpoints}
  \item IoU (Intersection over Union) measures overlap between prediction and ground truth. IoU = 1 is perfect, IoU = 0 is no overlap.
  \item Our IoU values are low (0.007--0.076) because kelp covers only a small fraction of each tile. This is \textbf{class imbalance}.
  \item Even with low absolute IoU, the \textbf{relative ranking} is meaningful.
\end{talkingpoints}

\subsection{Key Observations}

\begin{talkingpoints}
  \item \textbf{Oracle (0.0765)}: Best possible. Shows there's room for adaptive selection.
  \item \textbf{GNDVI (0.0699)}: Best single heuristic. Green-based index works well for kelp.
  \item \textbf{Ensemble (0.0695)}: Close second. Majority voting is robust.
  \item \textbf{NDVI (0.0533)}: Middle performer. Standard vegetation index.
  \item \textbf{Random (0.0445)}: Our floor. LA must beat this.
  \item \textbf{FAI (0.0074)}: Worst performer. Despite being designed for floating algae, it fails here---possibly due to threshold tuning or water conditions.
\end{talkingpoints}

\subsection{The Oracle Gap}

\begin{talkingpoints}
  \item Gap between Oracle (0.0765) and best Fixed (0.0699) is about 9\%.
  \item This gap represents the potential gain from adaptive selection.
  \item If LA can close this gap, it's learning something useful.
\end{talkingpoints}

%% ============================================================
\section{Slide 8: Week 2 -- Experiment 1 (Binary Reward, Failed)}
%% ============================================================

\subsection{Binary Reward Definition}

From \texttt{reward.py:126-129}:

\begin{lstlisting}[language=Python]
if self.reward_type == "binary":
    reward = 1.0 if iou > self.iou_threshold else 0.0
\end{lstlisting}

\begin{talkingpoints}
  \item Binary reward is a threshold-based signal:
  \begin{itemize}
    \item If IoU $> 0.3$ $\rightarrow$ reward = \textbf{1} (success)
    \item If IoU $\leq 0.3$ $\rightarrow$ reward = \textbf{0} (failure)
  \end{itemize}
  \item The threshold 0.3 was chosen arbitrarily---a ``reasonable'' value.
\end{talkingpoints}

\subsection{Why It Failed}

Your actual IoU values ranged from \textbf{0.007 to 0.076} (kelp is sparse). Since \textbf{no tile ever exceeded 0.3}, the automaton received:

\begin{center}
\begin{tabular}{lcc}
\toprule
\textbf{Tile} & \textbf{IoU} & \textbf{Reward} \\
\midrule
Tile 1 & 0.05 & 0 \\
Tile 2 & 0.07 & 0 \\
Tile 3 & 0.01 & 0 \\
\ldots & \ldots & 0 \\
\bottomrule
\end{tabular}
\end{center}

\begin{talkingpoints}
  \item Every heuristic got reward = 0, so the $L_{R-I}$ automaton had \textbf{no learning signal}.
  \item Probability updates became random drift.
  \item It converged to FAI \textbf{by chance}---the worst heuristic!
  \item This is a classic failure mode: \textbf{threshold miscalibration}.
\end{talkingpoints}

%% ============================================================
\section{Slide 9: Week 2 -- Fix with Continuous Reward}
%% ============================================================

\subsection{Continuous Reward Definition}

From \texttt{reward.py:128-129}:

\begin{lstlisting}[language=Python]
elif self.reward_type == "continuous":
    reward = iou  # Use IoU directly
\end{lstlisting}

\subsection{The Solution}

\begin{talkingpoints}
  \item Continuous reward uses IoU directly: $r = \text{IoU}$.
  \item Now every tile provides a gradient signal, even with low IoU values.
\end{talkingpoints}

Now the automaton sees meaningful differences:

\begin{center}
\begin{tabular}{lc}
\toprule
\textbf{Heuristic} & \textbf{Reward (avg)} \\
\midrule
FAI & $\approx 0.007$ \\
GNDVI & $\approx 0.07$ \\
\bottomrule
\end{tabular}
\end{center}

\begin{talkingpoints}
  \item Even though both are ``low,'' the \textbf{10$\times$ difference} provides a gradient signal.
  \item The automaton can now learn that GNDVI/Ensemble are better than FAI.
  \item The automaton can distinguish between ``bad'' and ``very bad.''
\end{talkingpoints}

\subsection{Why This Works}

\begin{talkingpoints}
  \item Binary reward is like pass/fail grading---you lose information.
  \item Continuous reward is like percentage grading---you preserve ordering.
  \item For problems with extreme class imbalance, continuous rewards are often necessary.
\end{talkingpoints}

%% ============================================================
\section{Slide 10: Week 2 -- Experiment 2 (Continuous, $\alpha=0.1$)}
%% ============================================================

\subsection{The Result}

\begin{talkingpoints}
  \item With continuous reward, LA correctly converged to Ensemble---one of the best heuristics!
  \item \textbf{But}: Convergence was too fast. By epoch 2, Ensemble had 99.9999\% probability.
  \item Entropy dropped to 0, meaning no exploration after epoch 2.
\end{talkingpoints}

\subsection{What is Entropy?}

Entropy measures how ``spread out'' the probability distribution over heuristics is.

\textbf{Formula:} For a probability distribution $p = [p_1, p_2, p_3, p_4]$ over 4 heuristics:
\[
H(p) = -\sum_{i=1}^{4} p_i \log(p_i)
\]

\textbf{What It Means:}

\begin{center}
\begin{tabular}{lcl}
\toprule
\textbf{Entropy} & \textbf{Distribution} & \textbf{Interpretation} \\
\midrule
High ($\approx$1.39) & [25\%, 25\%, 25\%, 25\%] & Uniform --- maximum exploration \\
Medium ($\approx$1.52) & [20\%, 0\%, 41\%, 39\%] & Some preference, still exploring \\
Low ($\approx$0.1) & [1\%, 0\%, 0\%, 99\%] & Strong preference, little exploration \\
Zero (0.0) & [0\%, 0\%, 0\%, 100\%] & Converged --- no exploration at all \\
\bottomrule
\end{tabular}
\end{center}

\textbf{In Your Experiments:}
\begin{talkingpoints}
  \item \textbf{Exp 2 ($\alpha$=0.1)}: Entropy dropped to \textbf{0.0} by epoch 2 --- the automaton ``locked in'' on Ensemble with 100\% probability. No more exploration.
  \item \textbf{Exp 3 ($\alpha$=0.01)}: Entropy stayed at \textbf{1.52} after 10 epochs --- GNDVI (41\%) and Ensemble (39\%) are still competing. The automaton is still exploring.
\end{talkingpoints}

\textbf{Why It Matters:}
\begin{talkingpoints}
  \item Low entropy too early = \textbf{premature convergence}. The automaton commits before gathering enough evidence.
  \item Higher entropy = continued learning and potential to find better solutions.
\end{talkingpoints}

\subsection{Why Fast Convergence is Bad}

\begin{talkingpoints}
  \item The automaton ``locked in'' before exploring alternatives.
  \item It converged to Ensemble, but GNDVI is actually slightly better (0.0699 vs 0.0695).
  \item With more exploration, it might have found GNDVI.
  \item Also, we can't trust results from only 1 epoch of real learning.
\end{talkingpoints}

%% ============================================================
\section{Slide 11: Week 2 -- Experiment 3 (Continuous, $\alpha=0.01$)}
%% ============================================================

\subsection{Proper Learning Dynamics}

\begin{talkingpoints}
  \item Learning rate $\alpha = 0.01$ (10x smaller than before).
  \item Now we see gradual convergence over 10 epochs.
  \item FAI is correctly eliminated (probability drops to ~0\%).
  \item GNDVI and Ensemble compete for dominance (41\% vs 39\% by epoch 10).
  \item Entropy = 1.52 means ongoing exploration.
\end{talkingpoints}

\subsection{Reading the Probability Table}

\begin{talkingpoints}
  \item Epoch 1: Roughly uniform (37\%, 3\%, 27\%, 34\%). FAI already penalized.
  \item Epoch 4: NDVI temporarily leads (44\%), FAI nearly gone (0.03\%).
  \item Epoch 7: Ensemble surges (53\%), competition continues.
  \item Epoch 10: GNDVI edges ahead (41\%), Ensemble close behind (39\%).
  \item This matches our baseline ranking: GNDVI $>$ Ensemble $>$ NDVI $\gg$ FAI.
\end{talkingpoints}

%% ============================================================
\section{Slide 12: Week 2 -- Experiment Comparison}
%% ============================================================

\subsection{Table Breakdown: What Each Row Means}

\begin{center}
\small
\begin{tabular}{lccl}
\toprule
\textbf{Metric} & \textbf{Exp 2} & \textbf{Exp 3} & \textbf{Explanation} \\
\midrule
Convergence & Epoch 2 & 10 epochs & How fast it ``decided'' \\
Final Entropy & 0.0 & 1.52 & 0 = no exploration; 1.52 = still exploring \\
Best Heuristic Prob & 100\% & 41\% & Exp 2 went ``all in''; Exp 3 hedged \\
Test IoU & 0.0719 & 0.0711 & Nearly identical performance \\
vs Random & +0.0151 & +0.0144 & Both $\sim$34\% better than random \\
Gap to Oracle & 6.0\% & 7.0\% & How far from perfect (0.0765) \\
\bottomrule
\end{tabular}
\end{center}

\subsection{Key Insight: Same Destination, Different Journey}

\begin{talkingpoints}
  \item Both experiments achieved similar final performance ($\sim$0.071 IoU).
  \item But Exp 2 (high $\alpha$) committed immediately, while Exp 3 (low $\alpha$) explored longer.
  \item Exp 2 got lucky---it happened to lock onto a good heuristic early.
  \item Exp 3 is more trustworthy---the gradual convergence shows the automaton actually \textbf{learned} the ranking.
\end{talkingpoints}

\subsection{Trade-offs}

\begin{center}
\begin{tabular}{lll}
\toprule
& \textbf{High $\alpha$ (0.1)} & \textbf{Low $\alpha$ (0.01)} \\
\midrule
\textbf{Pro} & Fast convergence & Robust, explores alternatives \\
\textbf{Con} & Risk of premature lock-in & Slower to converge \\
\bottomrule
\end{tabular}
\end{center}

\begin{talkingpoints}
  \item For \textbf{research}, prefer low $\alpha$ to observe learning dynamics.
  \item For \textbf{production} with time constraints, higher $\alpha$ may be acceptable.
\end{talkingpoints}

\subsection{Gap to Oracle}

\begin{talkingpoints}
  \item LA achieves ~94\% of Oracle performance.
  \item The remaining 6-7\% gap suggests room for improvement.
  \item Possible improvements: per-tile context, more heuristics, better reward shaping.
\end{talkingpoints}

%% ============================================================
\section{Slide 13: Week 2 -- Key Learnings}
%% ============================================================

\subsection{Technical Lessons}

\begin{talkingpoints}
  \item \textbf{Lesson 1}: Binary reward requires threshold calibration. Know your data distribution before setting thresholds.
  \item \textbf{Lesson 2}: Continuous reward is more robust for imbalanced problems.
  \item \textbf{Lesson 3}: Learning rate controls exploration/exploitation trade-off.
  \item \textbf{Lesson 4}: $L_{R-I}$ successfully learns heuristic quality rankings---the core hypothesis is validated!
\end{talkingpoints}

\subsection{Next Steps}

\begin{talkingpoints}
  \item Run more epochs (20-50) to see full convergence behavior.
  \item Try even smaller $\alpha$ (0.001) for more gradual learning curves.
  \item Add $\varepsilon$-greedy exploration: with probability $\varepsilon$, select randomly instead of following the learned distribution. This prevents probability collapse.
\end{talkingpoints}

\subsection{Broader Implications}

\begin{talkingpoints}
  \item The prototype validates the core thesis idea: LA can learn to select appropriate heuristics.
  \item Next phase: move from Global Policy (one distribution for all tiles) to Contextual Policy (different distributions based on tile features).
  \item Eventually: incorporate ML-based heuristics alongside spectral ones.
\end{talkingpoints}

%% ============================================================
\section{Slide 14: Next Steps -- Prototype Refinement}
%% ============================================================

\subsection{Prototype Refinement}

\begin{talkingpoints}
  \item Run more epochs (20--50) to see if distribution fully converges or keeps oscillating.
  \item Try $\alpha = 0.001$ for even more gradual learning curves.
  \item Add $\varepsilon$-greedy exploration: with probability $\varepsilon$, select randomly. Prevents probability collapse to 0/100\%.
\end{talkingpoints}

%% ============================================================
\section{Slide 15: Next Steps -- LA Variant Comparison}
%% ============================================================

\subsection{Available Algorithms}

All implemented in \texttt{automaton.py}:

\begin{center}
\small
\begin{tabular}{lll}
\toprule
\textbf{Algorithm} & \textbf{Description} & \textbf{Expected Behavior} \\
\midrule
L$_{R-I}$ & Reward only, ignore penalty & Conservative, stable \\
L$_{R-P}$ & Reward and penalty updates & Faster but noisier \\
Pursuit & Estimates reward, pursues best & 10--20x faster convergence \\
Discretized Pursuit & Finite probability levels & Even faster \\
Estimator (SERI) & Bayesian estimation + pursuit & Fastest $\varepsilon$-optimal \\
VSLA & Adaptive $\alpha$ based on entropy & Auto-balances explore/exploit \\
\bottomrule
\end{tabular}
\end{center}

\subsection{Talking Points}

\begin{talkingpoints}
  \item \textbf{L$_{R-I}$} (what we tested): Only updates on positive reward. Conservative---won't over-react to noise.
  \item \textbf{L$_{R-P}$}: Also penalizes bad actions. Faster learning but can oscillate in noisy environments.
  \item \textbf{Pursuit}: Maintains reward \textit{estimates} for each action, then ``pursues'' the best estimated action. Much faster convergence (10--20x) because it uses information more efficiently.
  \item \textbf{Discretized Pursuit}: Uses finite probability levels instead of continuous. Can converge even faster.
  \item \textbf{Estimator (SERI)}: Bayesian version of Pursuit. Maintains Beta distribution posteriors. Recognized as one of the fastest $\varepsilon$-optimal algorithms.
  \item \textbf{VSLA}: Automatically adjusts learning rate based on entropy. High entropy $\rightarrow$ high $\alpha$ (explore). Low entropy $\rightarrow$ low $\alpha$ (exploit).
\end{talkingpoints}

\subsection{Experiment Plan}

\begin{talkingpoints}
  \item Run all 6 algorithms on the same kelp detection task.
  \item Compare: convergence speed, final IoU, stability, entropy dynamics.
  \item Hypothesis: Pursuit/Estimator will converge faster but L$_{R-I}$ may be more stable.
\end{talkingpoints}

%% ============================================================
\section{Summary: What to Remember}
%% ============================================================

\begin{enumerate}[itemsep=0.5em]
  \item \textbf{The prototype works.} LA successfully learns to favor better heuristics.
  \item \textbf{Reward design matters.} Binary reward failed due to threshold miscalibration; continuous reward succeeded.
  \item \textbf{Hyperparameters matter.} Learning rate $\alpha$ controls convergence speed and exploration.
  \item \textbf{Baselines provide context.} Random (floor), Oracle (ceiling), Fixed (single-heuristic performance).
  \item \textbf{LA outperforms random selection.} This validates the adaptive selection approach.
\end{enumerate}

\end{document}
